\chapter{Введение}

\section{Цель}  

Произвести настройку маршрутизаторов.
Научиться настраивать протокол OSPF в рамках одной области на
маршрутизаторах для организации сетевой связности.
\section{Задачи}

\begin{enumerate}
\item Интерфейс Gi1/0/1 маршрутизатора R1 соединить с интерфейсом Gi1/0/2
маршрутизатора R2.
\item Интерфейс Gi1/0/2 маршрутизатора R1 соединить с интерфейсом Gi1/0/1
маршрутизатора R3.
\item Интерфейс Gi1/0/1 маршрутизатора R2 соединить с интерфейсом Gi1/0/2
маршрутизатора R3.
\item Компьютер VPC1 подключить к интерфейсу Gi1/0/3 маршрутизатора R1.
\item Компьютер VPC2 подключить к интерфейсу Gi1/0/3 маршрутизатора R2.
\item Компьютер VPC3 подключить к интерфейсу Gi1/0/3 маршрутизатора R3.
\item Настроить имена устройств и IP - адреса интерфейсов согласно таблице
5.2.1.
\item На маршрутизаторе R1 создать процесс OSPF 250, указать Router ID 1.1.1.1,
добавить маршрутизатор в область 0.0.0.0, активировать область, активировать
протокол.
\item На маршрутизаторе R2 создать процесс OSPF 250, указать Router ID 2.2.2.2,
добавить маршрутизатор в область 0.0.0.0, активировать область, активировать
протокол.
\item На маршрутизаторе R3 создать процесс OSPF 250, указать Router ID 3.3.3.3,
добавить маршрутизатор в область 0.0.0.0, активировать область, активировать
протокол.
\item На маршрутизаторах R1, R2, R3 привязать OSPF интерфейсы Gi1/0/1 и
Gi1/0/2, Gi1/0/3 к номеру процесса 250, указать какой области будут принадлежать
интерфейсы, интерфейс Gi1/0/3 сделать пассивным, активировать процесс OSPF
на интерфейсах.
\item Проверить таблицу маршрутизации.
\item Проверить сетевую связность VPC1, VPC2 и VPC1, VPC3 командами ping и
tracert.
\item Проверить сетевую связность между VPC1 и VPC2 командой ping с ключом -t
и в процессе выполнения, произвести разрыв соединения между
маршрутизаторами R1 и R2.
\item Проверить таблицу маршрутизации на маршрутизаторе R1.
\end{enumerate}

\begin{figure}[H] 
    \centering
    \includegraphics[width=0.75\textwidth]{top.png}
    \caption{Пример собранной сети}
\end{figure}

\begin{figure}[H] 
    \centering
    \includegraphics[width=0.75\textwidth]{table.png}
    \caption{Таблица с настройками}
\end{figure}

\endinput